\section{Folding}

Now the one real demand. To combine two distinctions is to multiply them. To combine them reversibly, losing no information, is to require that the multiplication never destroy anything. Made exact, this single demand pins down the entire algebra of the world.

Let the norm of an element be its length, and require the norm to multiply, so the length of a product is the product of the lengths. An algebra with such a norm is called a composition algebra, and from that one condition a chain of consequences follows. If the norm multiplies and two elements are nonzero, the length of their product is a positive number times a positive number, hence positive, so the product is nonzero. There are no zero divisors, no two nonzero elements whose product is zero. No zero divisors means every nonzero element has an inverse, so you can always divide, a division algebra. And division means reversibility: given a product and one factor, the other factor is recovered by multiplying by the inverse, with nothing lost.

So four statements are one and the same. The norm multiplies. There are no zero divisors. The algebra is a division algebra. The folding is reversible. To fold distinctions without losing information just is to live in a normed division algebra. This is the hinge of the whole construction, and it is derived: the informal demand of reversibility and the precise notion of a normed division algebra are the same requirement, not an analogy between them.

This is the one place the model asks something of the world beyond the bare premise, and it asks the least it can: only that combining differences should waste nothing. Everything specific about spacetime and matter follows from chasing that single requirement to its limit, which the next sections do.
